% New tenth-hour proof.  This is an actual ambient-retention theorem;
% it does not assert that the complete finite pentagon census has run.
\subsection{An actual nine-point finite quotient retaining the ordered-pole branch output}

Let $\mathcal M=\Gamma_{1,[9]}$ be the ordinary mixed mapping
class group at the configuration consisting of the origin and
four named opposite pairs.  Its named involution is
$\iota:z\mapsto-z$, with one fixed label and four two-element
label orbits.  Write
$\mathcal C=C_{\mathcal M}(\iota)$.
The actual branch quotient, forgetting the other three branch
points, and the monic ordered-pole double cover define the
ordinary homomorphism
\[
 \mathcal C\xrightarrow{\delta}
       \Gamma_{0,\{b_0;[t_1,t_2,t_3,t_4]\}}
       \xrightarrow{\operatorname{mon}}\Gamma_{1,2}.
                                                               \tag{FBO.1}
\]
Here the semicolon distinguishes $b_0$ and the brackets permit
permutation of the four root labels.  The two target elliptic
points are ordered.  The first arrow is the literal ordinary
map of (IFM.1)--(IFM.3); it retains the common permutation of
the other three branch points until those points are forgotten.
The monic equation and its two poles are those of (A.C6).

Use the actual quotient $\mathcal P_c=\Sigma/R_c$ of
(QF.1), where $\Sigma=F(A,B)$ and
$|\mathcal P_c|=32\cdot3^{14}$.  Its kernel is invariant under
the ordinary modular automorphisms.  The characteristic
source in (QF.3) is a separate larger quotient; no characteristic
invariance of $R_c$ under every automorphism of $\Sigma$ is
needed below.

\begin{result}{A.4.53}{Proved}{The final ordered-pole branch quotient is retained by an explicitly bounded characteristic surface representation of the original nine-point group}
There is a specified finite quotient $W$ of the two-punctured
torus group and an actual finite quotient
\[
 q_{\mathrm{amb}}:\mathcal M\twoheadrightarrow Q_{\mathrm{amb}}
                                                               \tag{FBO.2}
\]
with a well-defined homomorphism on the ordinary centralizer image
\[
 \theta:q_{\mathrm{amb}}(\mathcal C)
     \longrightarrow
     J_W:=\operatorname{im}\bigl(\Gamma_{1,2}
                            \to\operatorname{Out}(W)\bigr)
                                                               \tag{FBO.3}
\]
such that
\[
 \theta q_{\mathrm{amb}}|_{\mathcal C}
       =\rho_W\operatorname{mon}\delta.
                                                               \tag{FBO.4}
\]
The ordinary forgetting coordinate $J_W\to J_c$ and the
point-pushing homology coordinate modulo twelve are both
retained by their actual quotient maps.

Here are finite bounds specifying the construction.  Put
\[
 \begin{aligned}
 w&=12^3\,32\cdot3^{14},&d&=2w^2,&
 m&=d\,3^{1+3d},\\
 B_m&=(m!)^{(m!)^7},&
 v&=B_m3^{1+5B_m},&
 u&=(v!)^{(v!)^{10}}.
 \end{aligned}                                                   \tag{FBO.5}
\]
Then $|W|\leq w$ and $|Q_{\mathrm{amb}}|\leq9!u^{10}$.
Every displayed group is the image of a specified ordinary
homomorphism, so these are bounds on actual finite groups,
not bounds on independently chosen tuples.

Applying A.4.54 with $n=9$ to (FBO.2) gives an explicitly
bounded actual refinement $\eta:\mathcal M\to E$ whose
centralizer projects exactly to $q_{\mathrm{amb}}(\mathcal C)$.
Thus (FBO.3) also defines the literal homomorphism
\[
 C_E(\eta\iota)\longrightarrow J_W
                                                               \tag{FBO.6}
\]
by that projection.  This statement retains the final branch
output in an ambient finite group; it does not assert that a
transported action preserves an arbitrary quotient kernel.
\end{result}

\begin{proof}
\par\medskip\noindent\textbf{The two actual double covers.}\quad
Removing the origin and the eight paired points from the torus,
then taking its quotient by $\iota$, gives a sphere with the
distinguished puncture $b_0$, four punctures $t_j$, and three
order-two cone points $b_1,b_2,b_3$.  Its based orbifold group is
\[
 \mathcal E=
 \left\langle b_0,b_1,b_2,b_3,t_1,t_2,t_3,t_4\ \middle|\
 b_0b_1b_2b_3t_1t_2t_3t_4=1,\quad b_i^2=1\ (1\leq i\leq3)
 \right\rangle.
                                                               \tag{FBO.7}
\]
Eliminating $t_4$ gives
$\mathcal E=C_2(b_1)*C_2(b_2)*C_2(b_3)*
F(b_0,t_1,t_2,t_3)$.
The original branch-cover character is
\[
 \chi_B(b_i)=1\ (0\leq i\leq3),\qquad
 \chi_B(t_j)=0.
                                                               \tag{FBO.8}
\]
Its kernel $K=\Pi_{1,9}$ is the actual surface subgroup.
It meets each conjugate cone group trivially and is free.
The Euler characteristic of $\mathcal E$ is $-9/2$,
so $K$ has rank ten.

Forget the three cone points and impose $t_j^2=1$ on the four
remaining root meridians.  This is the actual quotient map
\[
 \mathcal E\twoheadrightarrow
 \Omega=C_2(t_1)*C_2(t_2)*C_2(t_3)*C_2(t_4),\qquad
 b_i\longmapsto1\ (i>0),\quad
 b_0\longmapsto(t_1t_2t_3t_4)^{-1}.
                                                               \tag{FBO.9}
\]
The new cover character is
$\chi_T(t_j)=1$, $\chi_T(b_i)=0$ for every $i$.
It defines the two-punctured torus subgroup
$K_*=\ker(\Omega\to C_2)$, of rank three.
In particular $\chi_B(b_0)=1$ and $\chi_T(b_0)=0$.
These are two different specified covering maps.

Here are complete words for the second surface marking.  Put
\[
 \begin{gathered}
 u=t_2t_1,\qquad v=t_3t_1,\qquad z=t_4t_1,\\
 A=u,\qquad B=u^{-1}v,\qquad
 C=u^{-1}vuv^{-1}zv^{-1}u,\\
 u=A,\qquad v=AB,\qquad z=[A,B]CB.
 \end{gathered}                                                   \tag{FBO.10}
\]
The first line is the Schreier free basis of $K_*$, and the last
line is the inverse to the middle change of free basis.  With
$t=t_1t_2t_3t_4$ one has the literal identities
\[
       t=([A,B]C)^{-1},\qquad
       C=B(t_1tt_1)B^{-1},\qquad b_0=t^{-1}.
                                                               \tag{FBO.11}
\]
Thus $t$ and $C$ are the two boundary orientations obtained by
removing the disk at infinity from the branched cover.  The
positive puncture meridians at its two poles are their inverses.
Filling the pole represented by $C$ is the based quotient
$\operatorname{cap}_C:K_*\to\Sigma=F(A,B)$ with $C\mapsto1$.
It is the specified pole-filling map, regardless of the chosen
orientation of that peripheral loop.

The ordinary Artin generators on $\Omega$ send
$(t_i,t_{i+1})$ to $(t_it_{i+1}t_i,t_i)$ and fix the other
generators.  They preserve $t$ as an element and preserve both
pole conjugacy classes separately.  Under $\operatorname{cap}_C$
their three actions are exactly
\[
 \begin{array}{c|cc|c}
 &A&B&\text{compact period}\\ \hline
 a&A&BA&U\\
 b&AB^{-1}&B&L\\
 c&A&BA&U
 \end{array}
 \qquad
 U=\begin{pmatrix}1&1\\0&1\end{pmatrix},\quad
 L=\begin{pmatrix}1&0\\-1&1\end{pmatrix}.
                                                               \tag{FBO.12}
\]
These are the signed folding maps of (A.C8), in a primitive
integral basis.  The earlier alternative basis $(v,u^{-1})$
is $(AB,A^{-1})$; this explicit commutator-preserving Nielsen
map supplies the comparison with that marking.

\par\medskip\noindent\textbf{The finite ordered-pole output.}\quad
Define $W$ to be the image of the actual joint map
\[
 K_*\longrightarrow(\mathbb Z/12)^3\times\mathcal P_c,
 \qquad g\longmapsto
       \bigl(g^{\mathrm{ab}}\bmod12,
                       \operatorname{cap}_C(g)\bmod R_c\bigr),
                                                               \tag{FBO.13}
\]
and let $N_*\triangleleft K_*$ be its kernel.  Formula
(FBO.12) and ordinary modular invariance of $R_c$ show that
every Artin generator preserves $N_*$.  This uses the actual
capped action, rather than invariance under an unspecified
automorphism.  The full four-strand twist acts on $\Omega$
as $\operatorname{Inn}(t)$; since $t\in K_*$, its restricted
action on $W$ is inner.  Consequently the outer action factors
through $B_4/\langle(abc)^4\rangle=\Gamma_{1,2}$, with the
monic ordered poles of (FBO.1).

For clarity, (FBO.13) retains the whole point-pushing homology
coordinate.  On $H_1(K_*;\mathbb Z)$ in the basis $(A,B,C)$,
the Artin matrices are
\[
 a_*=
 \begin{pmatrix}1&1&0\\0&1&0\\0&0&1\end{pmatrix},\quad
 b_*=
 \begin{pmatrix}1&0&0\\-1&1&0\\0&0&1\end{pmatrix},\quad
 c_*=
 \begin{pmatrix}1&1&0\\0&1&0\\0&-1&1\end{pmatrix}.
                                                               \tag{FBO.14}
\]
The two literal point-pushing generators of (A.C8),
$\xi_1=ca^{-1}$ and
$\xi_2=aba c^{-1}b^{-1}a^{-1}$, therefore have bottom rows
$(0,-1,1)$ and $(1,0,1)$, with identity compact blocks.
Their two homology coefficients are independent with coefficient
one.  In particular an inner action on $W$ has zero point-push
coordinate modulo twelve.  The projection $W\to\mathcal P_c$
also gives the actual forgetting quotient $J_W\to J_c$.
The ordinary outer image on $\mathcal P_c$ is exactly $J_c$
by A.4.22.  No homology coordinate is substituted for that
full finite ordinary image.

Now put
\[
 N=N_*\cap t_1N_*t_1^{-1},\qquad
 D=\Omega/N.
                                                               \tag{FBO.15}
\]
The subgroup $N_*$ is normal in $K_*$; conjugation by $t_1$
has order two and exchanges the two displayed factors.
Thus $N$ is normal in $\Omega$, and
$|D|\leq2|W|^2\leq d$.
Every based Artin automorphism $\gamma$ preserves $N$:
it preserves $N_*$, and $\gamma(t_1)t_1^{-1}\in K_*$,
whose inner action preserves both normal factors in (FBO.15).
The actual projections of $K_*/N$ to the two copies of $W$
keep the two pole coordinates separately.

\par\medskip\noindent\textbf{A graph layer detects the actual cusp conjugator.}\quad
Use the composite $q_D:\mathcal E\to D$ of (FBO.9),(FBO.15),
and put
\[
 L=\ker q_D,\qquad
 D^+=\mathcal E/(L'L^3).
                                                               \tag{FBO.16}
\]
The subgroup $L$ contains the three cone groups and is generally
not free.  Its graph of groups has all vertex groups of order
two.  Such a vertex group has zero positive homology over
$\mathbb F_3$: the averaging operator divides by two and
contracts its augmented resolution.  The graph-of-groups
homology sequence therefore identifies
$H_1(L;\mathbb F_3)$ with the first homology of its underlying
finite graph.

Use the graph of groups with one trivial central vertex,
three edges to the three $C_2$ vertices, and four oriented
free loop edges $b_0,t_1,t_2,t_3$ at the central vertex.
Because each cone generator maps to the identity in $D$,
the quotient graph for $L$ has $4|D|$ vertices and $7|D|$
edges.  Its first Betti number is $1+3|D|$.  Consequently
\[
                   |D^+|=|D|3^{1+3|D|}\leq m.
                                                               \tag{FBO.17}
\]
Let $r$ be the order of $q_D(b_0)$.  The homology class of
$b_0^r$ is the chain supported on the $r$ oriented $b_0$ edges
at vertices $1,q_D(b_0),\ldots,q_D(b_0)^{r-1}$, each with
coefficient one.  It is a nonzero cycle, also when $r=1$.
First homology of a graph injects into its edge-chain group.
Left deck translation stabilizes this chain exactly for the
elements of $\langle q_D(b_0)\rangle$.

An element of $D^+$ centralizing $b_0$ fixes this class under
conjugation.  Conversely each power of $q_D(b_0)$ has the
centralizing lift given by the same power of $b_0$.  We have
proved the exact finite projection
\[
 \operatorname{im}\bigl(C_{D^+}(b_0)\to D\bigr)
                         =\langle q_D(b_0)\rangle
                         \subseteq K_*/N.
                                                               \tag{FBO.18}
\]
The last inclusion uses $\chi_T(b_0)=0$, and is the pole-order
condition needed below.

\par\medskip\noindent\textbf{Reconstructing this output from a characteristic nine-point input.}\quad
For any finitely generated group $G$, let
$\kappa_s(G)$ be the intersection of kernels of all
homomorphisms $G\to\mathfrak S_s$.
Put
\[
 N_0=\kappa_m(\mathcal E),\quad P_0=\mathcal E/N_0,
 \quad N_1=N_0'N_0^3,\quad E_1=\mathcal E/N_1.
                                                               \tag{FBO.19}
\]
The left regular representation of $D^+$, enlarged to degree
$m$, shows $N_0\subseteq\ker(\mathcal E\to D^+)$.
Every character from $\mathcal E$ to $C_2$ is also retained
because $m\geq2$.  Thus $N_0\subset K$, and $N_0$ meets
every conjugate cone group trivially.  Its action on the
Bass--Serre tree is free, so it is a free group.  If
$p=|P_0|$, the actual graph has $p$ central vertices,
$3p/2$ cone vertices, and $7p$ edges.  Hence
\[
 \operatorname{rank}(N_0)=1+9p/2,
 \qquad |E_1|=p3^{1+9p/2}\leq v,
 \qquad p\leq B_m.
                                                               \tag{FBO.20}
\]
The bound on $p$ follows by listing all assignments of the
seven generators of $\mathcal E$ to $\mathfrak S_m$ and
retaining the assignments satisfying the three square relations.

The $P_0$ action on $H_1(N_0;\mathbb F_3)$ is faithful.
Here are both steps, including the passage from integral
homology to characteristic three.  A nonidentity deck element
fixes no oriented edge: every original edge group is trivial.
It fixes no central vertex and may fix some cone vertices.
The exact graph chain sequence gives its rational trace on
$H_1$ as $1-f$, where $f\geq0$ is its number of fixed cone
vertices.  The identity trace is $1+9p/2>1$.  Thus every
nonidentity deck element acts nontrivially on integral homology.
These integral matrices have finite order.

A finite-order integral matrix congruent to the identity modulo
three is the identity.  Indeed over $\mathbb Z_3$ write a
nonidentity such matrix as $I+3^aB$, with $a\geq1$ and
$B\not\equiv0\pmod3$.  Raising to a positive exponent prime
to three leaves that valuation unchanged, and the binomial
expansion gives
\[
 (I+3^aB)^3=I+3^{a+1}B\pmod {3^{a+2}}.
                                                               \tag{FBO.21}
\]
Successive cubings increase the finite valuation by exactly one;
no positive power is the identity.  This contradicts finite
order and proves the assertion.  Applying it to the deck
matrices proves the claimed faithful mod-three action.
In particular, if $K_1=K/N_1$, then
\[
                   C_{E_1}(K_1)\longrightarrow P_0
                       \quad\hbox{has image }1,
                                                               \tag{FBO.22}
\]
because such a centralizer acts trivially on
$N_0/N_1\leq K_1$.

Now define the characteristic surface quotient and its actual
index-two extension by
\[
 M=\kappa_v(K),\qquad Q_*=K/M,
 \qquad E_2=\mathcal E/M.
                                                               \tag{FBO.23}
\]
The regular representation of $K_1$, enlarged to degree $v$,
gives $M\subseteq N_1$.
The subgroup $M$ is characteristic in the free rank-ten group
$K$.  Since $K\triangleleft\mathcal E$, it is normal in
$\mathcal E$; characteristicness of $K$ in $\mathcal E$ is
not required.  Therefore all maps
$E_2\to E_1\to P_0\to D^+\to D$ are literal quotient maps,
and $|Q_*|\leq u$.

Set $Z=C_{E_2}(Q_*)$ and let $\mathcal A$ be the actual
image of $E_2$ in $\operatorname{Aut}(Q_*)$.
Equation (FBO.22) gives
\[
 Z\subseteq N_0/M\subseteq Q_*,\qquad Z=Z(Q_*),
 \qquad\mathcal A=E_2/Z.
                                                               \tag{FBO.24}
\]
The induced deck $C_2$ therefore embeds in
$\operatorname{Out}(Q_*)$, and $\mathcal A$ is its entire
inverse image in $\operatorname{Aut}(Q_*)$.
Every map $\mathcal E\to\mathfrak S_m$ kills $N_0$ and
factors through $\mathcal A$ by (FBO.24); conversely each map
from $\mathcal A$ gives such a map from $\mathcal E$.
Consequently
\[
 \kappa_m(\mathcal A)=(N_0/M)/Z,
 \qquad\mathcal A/\kappa_m(\mathcal A)=P_0.
                                                               \tag{FBO.25}
\]
An automorphism of $Q_*$ normalizing the named outer deck
group conjugates $\mathcal A$ to itself and thus descends to
the actual intrinsic quotient $P_0$ in (FBO.25).
Changing it by an inner automorphism of $Q_*$ changes the
descended automorphism by the corresponding inner automorphism
of $P_0$.  This is the actual finite inner-action reconstruction;
it has not assumed the vanishing of an extension-class obstruction
at the upper extension $E_2$.

\par\medskip\noindent\textbf{The kernel containment proving the ambient map.}\quad
Since $M$ is characteristic in $K$, the ordinary mapping-class
action gives $\mathcal M\to\operatorname{Out}(Q_*)$.
Join it with the actual nine-label permutation and define
$q_{\mathrm{amb}}$ as their joint image.  A group generated
by ten elements has at most $|Q_*|^{10}$ automorphisms, so
the bound $9!u^{10}$ follows.

Take $g\in\mathcal C$ with $q_{\mathrm{amb}}(g)=1$.
Ordinary branched-cover descent identifies its action with
an outer automorphism of the original orbifold group
$\mathcal E$.  A based disk representative of its sphere
mapping class fixes $b_0$ as an element: remove a small disk
at the distinguished branch point and use its boundary
basepoint.  Filling that disk changes representatives by
a boundary twist, whose group action is an inner power of
this same $b_0$.

There are two lifts of that based boundary representative to
the original double cover.  Their outer surface classes are
$g$ and $g\iota$.  Since $q_{\mathrm{amb}}(g)=1$, a representative
of either lift on $Q_*$ lies in $\mathcal A$, the full inverse
image of the named outer deck group in (FBO.24).
Its conjugation action on $\mathcal A$, and consequently on
the intrinsic quotient $P_0$ of (FBO.25), is inner.
Thus the chosen based orbifold action is inner on $P_0$,
say by $z\in P_0$; no choice between the two surface lifts
is suppressed here.
It fixes $b_0$ as an element, and hence $z$ centralizes $b_0$
in $P_0$.  Projecting through the retained $D^+$ and using
(FBO.18) gives
\[
                       z_D\in\langle q_D(b_0)\rangle
                                  \subseteq K_*/N.
                                                               \tag{FBO.26}
\]
All based ordinary sphere actions preserve the quotient to $D$:
they first forget the three cone points, then induce the actual
four-strand Artin action already checked in (FBO.12)--(FBO.15).
Thus their induced action on $D$ is the inner action by $z_D$.
Because $z_D$ is in the even subgroup, its restriction to
$K_*/N$ and then to $W$ is inner.  This proves
\[
 \ker(q_{\mathrm{amb}}|_{\mathcal C})
          \subseteq\ker(\rho_W\operatorname{mon}\delta).
                                                               \tag{FBO.27}
\]
Kernel containment for these two actual homomorphisms defines
the unique map (FBO.3), and gives (FBO.4).
The evenness assertion is essential: it retains the ordered
pole instead of treating an odd orbifold conjugator as an
inner surface conjugator.

Finally A.4.54 applies to the actual quotient (FBO.2), with
its prescribed label and relative-homology joins.  Its equality
$C_E(\eta\iota)\to q_{\mathrm{amb}}(\mathcal C)$ proves
(FBO.6).  This uses that proved ambient recognition theorem
after the exact kernel containment, rather than postulating
extension of an arbitrary finite quotient of $\mathcal C$.
\end{proof}

The executed certificate
\path{work/tower_code/finite_branch_surface_source.py}, with
its JSON and text output, checks the two free-basis maps,
all Artin inverse and braid identities, the two pole words,
the capped actions (FBO.12), every entry of (FBO.14), both
point-pushing rows, and the full twist as the inner action
of the even boundary.  The graph and characteristic-quotient
construction is proved above with explicit finite bounds;
its enormous generated groups have not been exhaustively
enumerated.

\par\medskip\noindent\textbf{Local same-parameter continuation.}\quad
The diagonal-disk calculation (DDP.1)--(DDP.11) gives the
literal local parameter record for the final output in
(FBO.4).  It is made on the original nine-point torus, so
no unproved preservation of an invariant pants decomposition
after adding three punctures is required.  Its ordinary
branch map is the same (FBO.1), through (IFM.1).
The remaining global finite compiler must place the translation
fusion, the involution recognition, and this retained output
in their common characteristic tower level.  No completed
finite gentle-shadow census or final parity exclusion is
asserted by the present retention theorem alone.
